\section{Namen Gottes}

Ich möchte heute diese Serie über die Namen Gottes abschliessen. Zwölf Sonntage habe wir uns jetzt mit den Namen Gottes beschäftigt. Beim recherchieren über die Namen, habe ich sehr viel von Gott gelernt und ich habe ein anderes Bild vom ihm bekommen. Er wurde mir noch grösser und ich verstehe ihn noch weniger.

Wenn wir uns mit den Namen Gottes beschäftigen, sollten wir das immer bewusst mit einer Gewissen Ehrfurcht tun weil Gott uns in seinem zweiten Gebot explicit darauf aufmerksam macht \biblezit{IIMos}{20:7}.

Wir haben mit den drei Hauptnamen Gottes angefangen.
\begin{itemize}
    \item \textbf{Elohim} Genereller Name für einen Gott
    \item \textbf{Adonai} Meister oder Herr
    \item \textbf{Jahwe} JHWH ist der Name, mit dem sich Gott Moses vorgestellt hat \biblezit{IIMos}{3:14}. 
\end{itemize}

Aus diesen drei Namen wurden dann diverse Namen zusammengesetzt, die das Wesen und die Eigenschaft unseres Gottes abbilden. Alle Namen, die wir angeschaut haben, finden wir so in der Bibel. Manche wurden von Menschen Gott gegeben wie in, \biblezit{IMos}{16:13}. Andere Namen kommen von Gott selber, wie zum Beispiel der Name \betonung[b, p]{Herr der Heerscharen} der Gott in Jeremia sich selber gegeben hat

Gott einen Namen zu geben ist immer ein Versuch, Gott verstehen zu wollen. Gott zu verstehen ist wie ein Kaleidoskop. Früher als Kinder hatten wir solche Röhrchen, wo man durchblicken konnte und man sah immer wieder eine neue Anordnung von Farben und Formen. So ist es auch mit den Namen Gottes. Sie wiederspiegeln immer eine Momentaufname von Gott. Einmal ist er ein eifersüchtiger Gott, du drehst dich um und dann ist er wieder ein gerechter Gott. 

Aber es ist komplizierter. Gott ist nicht einmal der Arzt und ein anderes Mal der Richter. Er ist immer alles zur gleichen Zeit. Gott ist jetzt hier bei uns, ist aber überall zur gleichen Zeit. So ist sein Wesen. Er ist immer alles zur gleichen Zeit. Ich begreife das nicht. Ich weiss aber, dass ich mich auf ihn verlassen kann. Auch wenn es nicht immer so wird wie ich es mir wünsche oder gewünscht habe, Gott macht aber keine Fehler. Wir werden eine Ewigkeit Zeit haben, um zu versuchen, zu verstehen, wer Gott ist. Der Vers \biblezit{IPetr}{1:12} sagt aber, dass nicht einmal die Engel das Wirken Gottes total verstehen. Der Satan weiss es schon gar nicht, wie wir das in \biblezit{Mk}{1:24} lesen.

Darum sagt Gott, wir sollen kein Bild von ihm machen. Egal wie wir ihn uns vorstellen, er ist anders. Diese Serie hat mir gezeigt, wie gross Gott ist und wie wenig ich von ihm fassen kann. Und wisst ihr was? Das ist genau das, was mich tröstet! Was wäre das für ein mikriger Gott, wenn ich ihn erklären könnte? 

Wir dürfen diesen Gott anbeten, über ihn staunen und uns an ihm erfreuen. Und durch den Herrn Jesus Christus ist nun der Weg offen zu ihm. Dafür bin ich dankbar.

So schliesse ich diese Serie mit einem Spruch aus einem jüdischen Gebetsbuch:
\begin{quote}
    In Gottes Gebot ruht unendlichte Macht,\\
    die Worte nicht beschreiben können.\\
    Wären alle Himmel Pergament,\\
    alle Schilfrohre Schreibfedern, alle Meere Tinte\\
    und alle Bewohner der Erde Schreiber,\\
    so könnte Gottes Herrlichkeit doch nicht vollständig erzählt werden.
\end{quote}

